\documentclass[journal]{IEEEtran}
\usepackage[hidelinks]{hyperref}
\usepackage{booktabs}
\usepackage{enumitem}
\usepackage{microtype}
\usepackage{float}
\usepackage{array}
\setlength{\emergencystretch}{1em}

\title{When Abstraction Becomes Indirection:\\Tokens-to-trace, measured on six stacks}
\author{Vasili Gavrilov}
\date{August 2026}

\begin{document}

\twocolumn[{%
\begin{center}
{\Huge\bfseries When Abstraction Becomes Indirection\par}
\vspace{6pt}
{\LARGE Tokens-to-trace, measured on six stacks\par}
\vspace{12pt}
{\large Vasili Gavrilov\par}
\end{center}
\vspace{10pt}
{\small\noindent\textbf{\textit{Abstract}:} Programs are full of layers made for human readers: frameworks, class
hierarchies, syntactic sugar. We ask whether code written with AI
agents still needs them. Before an agent can change code safely it
must read everything that decides what the code does. We count that
reading as tokens-to-trace: for one entry point, the tokens of the
smallest set of source pieces that decide its behavior, and where
each piece lives: in the project, in library source, or only in the
version's documentation. One path per stack, three pairs. Plain C
against idiomatic C++, whole files: 90 against 257 thousand tokens,
and held to the same behavior the C++ text is only 15\% bigger, so
the gap is scope and file scatter, not syntax. Plain Java with JDBC
against Spring: 7.6 thousand, all in the project, against 8.2
thousand in the project plus 51 to 124 thousand inside the framework
jars, and part of what Spring does is written in no file at all.
Vanilla JS against React: the same 1.1 thousand to write, plus 283
thousand of React's own machinery outside the project. Then 430
agent runs on the same pairs. C against C++: the C++ side cost about
20\% more turns, and a control pair shows the agent pays for the
walk across files, not for the dispatch. Java: one endpoint built on
both stacks and run in two batches of 40; the turn cost reversed its
sign between the batches, and what replicated is the failure: five
silent wrong answers, all Spring's, each a wire name derived from a
Java identifier by a convention no source file states. The web: a
change that crosses a component boundary cost React 100\% more turns
than vanilla. What a
reader must know first is also not equal: K\&R is 272 pages, the
Stroustrup 1,368, and a framework's documentation is versioned:
what it says and how much of it there is depend on the version.
One line in the prompt, ``plain C'', ``plain Java with JDBC'' or
``plain vanilla JavaScript'', gives the agent the whole truth of a path for a few thousand tokens;
the layered styles fill its memory, or their deciding text is not in
the project at all.\par}
\vspace{22pt}
}]




\section{The question and the measure}
\label{sec:measure}

A class hides a representation and an interface hides an implementation.
A framework hides something larger: control flow. Each is a remedy for a limit of human working
memory. An agent is not under that limit and is under others: the context it
carries and the relationships it must resolve before acting. A hidden
implementation is a relationship to resolve, which abstraction adds rather
than removes.

The machine-cost objection to abstraction was answered decades ago by
compiler work and the zero-overhead principle, and nothing in this paper
rests on machine cost (the wall-time numbers are in the
five-language bench in \texttt{ais} \cite{ais}). The reader's cost was never
priced, and the reader has changed since.

$\mathrm{tokens\mbox{-}to\mbox{-}trace}(e)$ is the token count of the minimal
set of source regions that fully determines the behavior of one entry point
$e$, wire to storage and back: routing, authentication, authorization, handler
logic, query, schema, serialization, error mapping. A region is in the set iff
removing it leaves some behavior of $e$ undetermined. The platform floor (JDK,
libc, servlet container, database engine, vendored third-party crypto) is
excluded identically everywhere; anything the application configures, or that
varies by library version, is included. Each region is classified by
residence:

\begin{enumerate}[nosep]
\item \textbf{repository}: read directly from the project's own source;
\item \textbf{library source}: text outside the repository that decides
behavior before the program runs;
\item \textbf{convention}: fixed by defaults of the exact version in use,
recoverable from its documentation or from a runtime observation, not from
any source file.
\end{enumerate}

In the agent's own terms the closure is the needed context: what must be
in front of it before it can decide that a change is safe. Template
specialization and virtual-override resolution are classified by the
residence of the deciding code. Tokens are o200k, the encoding of the current OpenAI models, counted with
tiktoken, OpenAI's open-source tokenizer library \cite{tiktoken}: a public,
model-neutral count that approximates what a commercial agent is billed (the
pilot's own billing used the model vendor's tokenizer, which differs); the region manifests, with per-row lines, characters and
tokens, are \texttt{closures.py} and \texttt{closures.json} beside this
paper. The measure is a count: anyone can recompute it from the manifest, so two
systems compare by number, and a disagreement is about a manifest row, not
about an interpretation.

\section{Four closures}
\label{sec:closures}

One read path per system, each doing the same kind of work so the numbers
compare: an authenticated or filtered list over a local store, returning
rows.

\textbf{ais, plain C.} An associative index \cite{ais}, 10,597 code lines,
written to an explicit style guide \cite{aisstyle} in the school of K\&R
and Robbins \cite{knr,robbins} (MISRA-inspired; compliance with MISRA C:2012
\cite{misra} is not claimed or checked). The non-interactive recall path
(\texttt{ais KEY1 KEY2}: AND-intersection over posting lists, record
printing, document-blob branch) closes in \textbf{489 lines across 6 files,
5,750 tokens} with the storage engine as floor; for the Java pair below, MySQL is
that floor. The interactive secret-reveal branch, taken only at a tty,
adds 3,634 tokens including the contract of the crypto wrapper (the
vendored Monocypher beneath it counts as floor, a rule added in this
revision with Monocypher its only instance). With the hand-written engine
added (posting merge, key folding, ledger scan, tombstones, crash
recovery, headers as contracts), the determination of the non-interactive path from
\texttt{argv} to the bytes on disk is \textbf{1,345 lines across 16 files,
14,444 tokens}: 7\% of one window. The rows overlap by one 207-token
comment (5,750 + 8,901 - 207 = 14,444); with the reveal branch the total is
18,078, and then nothing but Monocypher is floor. Whole files over the same
16: 90,151. The callbacks on the path are function pointers passed at the
call site; the ops-struct vtable idiom of sqlite or nginx would resolve less
cheaply, so this row measures one explicit style, and an abstract-C row
(nginx's request path) is future work.

\textbf{Monocypher 4.0.3, third-party plain C.} To bound how much of the
plain rows is the author rather than the style, the same rule was applied
to a library nobody here wrote: Monocypher \cite{monocypher}, the cipher
ais vendors, widely audited, and the paper's own example of the
algorithmic style, functions over buffers. Its whole AEAD encrypt path
(\texttt{crypto\_aead\_lock}: the XChaCha20 keystream, the Poly1305
tag, the wipe) closes in \textbf{715 lines across 2 files, 7,560
tokens}, beside System~A's 7,566 and ais's 5,750. The posture differs by
design, a cipher rather than a list, and that is the point: at the
algorithmic extreme the plain style's closure is the same few thousand
tokens whoever writes it.

\textbf{taskwarrior 2.6.2, idiomatic C++.} A widely used open-source
command-line task manager, in development since 2006 \cite{taskwarrior},
whose \texttt{task list} does the same kind of work: filter records from a local
plain-text store, print a table. Version 2.6.2 is the last all-C++ release
(3.x moved the core to Rust), which is the reason it is the subject. Its two
dispatch registries, commands and columns, are filled once at startup from
literal constructor calls, so they resolve statically at their registration
sites, one hop more than a call by name; an earlier draft of this paper
called that granularity undefinable, which was wrong, and the retraction
corrects the claim, not the number, since the measured set
(\texttt{CmdCustom}, the 13 columns of the default report, their five base
classes) was already the resolved one. The closure at whole-file
granularity is \textbf{99
files, 33,706 lines, 256,664 tokens}, with 43 \texttt{virtual} declarations
and 26 \texttt{override}s on the path; with the store (its TDB2 class and file layer) as
floor, 239,699. Compared at the same granularity, whole files on
both sides, the C closure is \textbf{90k against the C++ 257k, a factor of
2.8 in tokens and 6 in files (16 against 99)}; regions inside files were not
cut out on either side of that comparison, so both numbers are upper bounds
of the same kind. A single
system does not establish the idiom's cost: the whole of leveldb, a well-factored idiomatic C++
codebase, is 136k o200k tokens, so its read-path closure would land far below
taskwarrior's; n here is 1, and the feature scopes of \texttt{task list} and
\texttt{ais} recall differ (urgency, color rules, virtual tags on one side).

\textbf{System~A, plain Java with JDBC.} A commercial web application in
production, 39,677 code lines, reported anonymously. The endpoint (token
authentication, role gate, per-object ACL, a 4-table JOIN, streamed JSON,
central error mapping) closes in \textbf{813 lines across 17 files, 7,566
tokens}: 3.8\% of a 200k context window. The schema enters as the frozen
base DDL plus the forward-only migrations that alter the seven tables the
query and its ACL clause touch, because that sum is the live schema: the
base is never rebuilt and every change moves forward, so the data survives
every revision. Reading every touched file whole instead of the exact
regions costs roughly 100k, still one window, and every callee on the path
is named literally at its call site, so each hop is findable by one search.
Categories 2 and 3 measure zero at the stated floor.

\begin{figure*}[!t]
\footnotesize
\textbf{Side note: the same parser, written plain.} The filter machinery on
the measured path is a small language: \texttt{Lexer}, \texttt{Variant}
with its 30 operator overloads, and \texttt{Eval}, 4,934 lines across 6
files. The prompt that asks for its plain form is one line: ``plain C,
ais-style, MISRA-inspired, the way ciphers and codecs are written: structs
and functions, fixed bounds.'' What that form looks like:

\begin{verbatim}
struct token  { enum { T_NUM, T_STR, T_DATE, T_OP, T_END } kind;
                const char *s; size_t n; };

struct value  { enum { V_NUM, V_STR, V_DATE, V_BOOL } kind;
                union { double num;
                        struct { const char *s; size_t n; } str;
                        long   date;
                        int    flag; } u; };

struct cursor { const char *p; };

static int next_token(struct cursor *c, struct token *t);
        /* one switch over *c->p; no allocation, the token
           points into the input */
static int eval(struct cursor *c, const struct task *t,
                struct value *out);
        /* recursive descent; each grammar rule one function,
           each operator one case of one switch */
\end{verbatim}

The 30 overloads become cases of the switch in \texttt{eval}; the three
classes become two structs and a cursor; the token and the value live on
the stack and point into the input, so the memory cost is the input plus a
few dozen bytes. This is a sketch, not a measured artifact: the fair
experiment would be to write it, pass taskwarrior's own tests, and count
it,
and the pairs of Table~\ref{tab:pairs} already price the constructs it
removes. It is shown because the claim ``the plain form exists and is
small'' is otherwise abstract.
\end{figure*}

\textbf{Spring Boot 4.1.0, framework Java.} \texttt{raeperd/}\allowbreak\texttt{realworld-}\allowbreak\texttt{springboot-java}
v2.1.1 \cite{realworld}, an implementation of RealWorld, the community
specification of one Medium-style blogging application that many stacks
implement identically so the implementations can be compared; this one is
conventional, well-regarded Spring (Framework 7.0.8, Security 7.1.0, Data
JPA 4.1.0, Hibernate 7.4.1, Jackson 3.1.4, resolved from the Boot BOM). Endpoint
\texttt{GET /articles?author=...}: authenticated, paginated, viewer-dependent
flag, five entities. Category 1: \textbf{27 files, 1,267 lines, 8,175
tokens}, 8\% more than System~A's whole closure for a domain without study
design, supervision or reporting behind it. Category 2, the 27 framework
classes on the dispatch path, measured from the published sources jars of the
exact versions:

\begin{table}[H]
\centering
\footnotesize\setlength{\tabcolsep}{4pt}
\begin{tabular}{@{}lrrr@{}}
\toprule
Module & Classes & Lines & o200k \\
\midrule
spring-webmvc & 9 & 6{,}310 & 52{,}178 \\
spring-web & 4 & 1{,}310 & 10{,}758 \\
spring-security (web, config, core) & 6 & 3{,}348 & 29{,}498 \\
spring-data-commons & 5 & 1{,}959 & 14{,}572 \\
spring-data-jpa & 3 & 2{,}282 & 16{,}788 \\
\midrule
Total & 27 & 15{,}209 & 123{,}794 \\
\bottomrule
\end{tabular}
\caption{Category 2: framework source whose text decides the endpoint's
behavior, whole classes on the dispatch path. The class list is one measurer's trace of the
path; a blind second rater's trace of the same path measures 373k tokens
(the threats section reports the agreement), so this table is the
conservative trace. Beneath both, uncounted in any closure,
hibernate-core measures 6.7M tokens of source (one tiktoken pass,
recorded in \texttt{closures.json}).
Drawing the floor above MVC dispatch and the filter chain, on the argument
that they are as stable as a servlet container, leaves 51,342: the
interpreters of this application's own text (derived-query parsing,
\texttt{Pageable} resolution, the repository proxy, and the
\texttt{HttpSecurity} DSL that turns its security configuration into a
filter chain).}
\label{tab:cat2}
\end{table}

Category 2 says where the deciding text lives, not what an agent reads: in
practice agents finish Spring tasks without opening framework source at all,
because the conventions sit in the training weights (that is practitioner
report, not a measurement of this paper), and this moves the cost into
category 3, where the corpus is a superposition of versions, rather than
removing it. The failures practitioners report agents making on Spring, again
report rather than measurement, are category-3-shaped: the \texttt{javax}-to-\texttt{jakarta}
migration, the removal of \texttt{WebSecurityConfigurerAdapter} in
Security~6, \texttt{@Transactional} self-invocation not proxying, lazy
loading outside a session, dialect drift between H2 and production. The
plain stack ate the \texttt{jakarta} rename too; only the surface differs: one import
namespace with identical semantics against a DSL whose semantics moved. Category 3 for this endpoint, enumerated: the derived-query grammar
(\texttt{findAllByAuthorProfileUserName} becomes a join tree by property-path
rules; the SQL is in no repository file and is recovered by turning on
Hibernate's statement logging for one run, or reconstructed from the rules);
mapping precedence among four same-URL methods discriminated by
\texttt{params}, an idiosyncrasy of this application (the common idiom, one
method with optional parameters, removes that machinery from the path);
\texttt{Pageable} defaults; which beans auto-configuration creates;
filter-chain order from the \texttt{HttpSecurity} DSL; JPA fetch and flush
defaults; serialization defaults; H2's PostgreSQL emulation. The Jackson~3
package rename (\texttt{tools.jackson}) and the Boot~4 module split are the
version-incoherence instances observed during measurement.

\begin{figure*}[!t]
\footnotesize
\begin{verbatim}
plain C              idiomatic C++        plain Java           framework Java
90k by whole files   257k by whole files  7.6k, all in repo    8.2k + 51k-124k out

main                 main                 dispatcher servlet   {filter chain} auth
  |                    |                    |                    |
arg parse + dispatch app object run       route + auth         framework dispatcher
  |                    |                    |                    |
query                {command registry}   endpoint branch      {route mapping}
  |                    |                    |                    |
store read           command object       DAO query            {argument binding}
  |                    |                    |                    |
print records        query                SQL read             controller + service
                       |                    |                    |
                     {operator overloads} rows to JSON         {tx proxy}
                       |                                         |
                     store read                                repository iface
                       |                                         |
                     {formatter factory}                       {query derivation}
                       |                                         |
                     print records                             {dynamic SQL} read
                                                                 |
                                                               {serializer} to JSON

  |   a step to a name at the call site, resolvable by one search
 {}  a step the reader cannot follow by the name at the call site: a
     registry, an interpreter, an overload set, a derivation: in-repo
     for the C++, in framework source or convention for the framework Java
\end{verbatim}
\caption{The path a reading agent resolves through the four measured
endpoints, the plain and layered member of each language pair side by
side. Every step is a role of the style, not a name from the project, so
the columns contrast the styles: each layered column is its plain
sibling's path with the style's resolution steps inserted. In the
measured C++ the registry is the application object's command map, the
command object one class per command, the overload set the value type's
30 operators, the factory the output-column registry; in the measured
framework the dispatcher is Spring's \texttt{DispatcherServlet}, the
route mapping its annotation conditions, the argument binding its
parameter and paging resolvers, the query derivation \texttt{PartTree}
over the repository method name, the serializer Jackson. The columns are
matched in the kind of work, an authenticated or filtered list over a
local store, not in feature scope; on the framework path the SQL and the
serialization defaults have no file of their own in the repository, and
the closures under the headings differ by the factors of
Table~\ref{tab:quad}. The inserted steps are not drawing choices: the
pairs experiment prices them (Table~\ref{tab:pairs}), and its control
shows the one-class-per-file layout, not the dispatch construct, carries
most of the largest cost.}
\label{fig:paths}
\end{figure*}

\begin{table}[H]
\centering
\footnotesize\setlength{\tabcolsep}{1.3pt}
\begin{tabular}{@{}llrrr@{}}
\toprule
System & Style & Cat.~1 & Cat.~2 & Cat.~3 \\
\midrule
ais, engine as floor & plain C & 5{,}750 & 0 & 0 \\
ais, engine included & plain C & 14{,}444 & 0 & 0 \\
taskwarrior, whole-file$^{*}$ & idiomatic C++ & 256{,}664 & 0 & 0 \\
System~A & plain Java+JDBC & 7{,}566 & 0 & 0 \\
Monocypher AEAD (third party) & plain C & 7{,}560 & 0 & 0 \\
Spring endpoint & framework Java & 8{,}175 & 123{,}794$^{\dagger}$ & $>$0$^{\S}$ \\
items screen twin & plain JS & 1{,}092 & 0 & 0 \\
items screen twin & React 19 & 1{,}145 & 282{,}717 & $>$0 \\
\bottomrule
\end{tabular}
\caption{Tokens-to-trace by residence, one read path per system, o200k.
Category 1 is text in the system's own repository, category 2 text in
library or framework source outside it, category 3 behavior fixed by the
version's conventions, with no source text anywhere. A zero means every
region that decides the path's behavior is in the repository: the first four rows use libraries only below the floor, and
their dispatch machinery, taskwarrior's included, is their own code, while
the framework row's deciding text sits partly in the Spring jars and
partly in version convention.
$^{*}$Upper bound; measured the same way, whole files, the C figure is 90k, factor 2.8.
$^{\dagger}$51k under the highest defensible floor (Table~\ref{tab:cat2}).
$^{\S}$Enumerated in the text; for the two-stack endpoint of
Section~\ref{sec:endpoint} the category is counted, 31,535 tokens of
documentation at the exact versions.}
\label{tab:quad}
\end{table}

Information
hiding lets a client not read the implementation: a well-typed client
cannot observe it, so for a client-side task the interface and its
contract are the whole read, and the closure measured here is an upper
bound the reader rarely pays. The runs agree: the median run read three
or four files, not the closure. The closure matters when that protection
stops working: a defect or a semantics change sits behind the boundary
and the reader must go in. The cross-module tasks of the pairs
experiment simulate that case, and maintenance eventually produces it. The residence claim is
untouched by the objection either way: a convention that is in no
repository cannot be read at any granularity. The next measurement this distinction asks for is task-conditioned
reading beside worst-case closure, on the same systems.

Each kind of step has a concrete cost to the reader, set by the
resolving procedure, and the experiment's numbers are consistent with it
at the ends. A name at the call site is one search: the callee's text is one grep
away, which is why the plain runs are a few reads and an edit. A registry
adds one search before the first one can land, finding the registration
site; measured, its pair's cost is mostly the walk across the idiom's
files, and what is left for the dispatch itself does not resolve against
noise. A template call resolves in one search once the argument types are
known, $+$1.70 measured. An overloaded operator leaves a greppable token
but not a unique name: the reader must establish the operand's type
before the grep lands, a cost too small to resolve in these pairs, and
argument-dependent lookup, which can put the deciding candidate in any
associated namespace under a rule written in no source file, is the
harder form of the same cost, not exercised by these pairs. Inheritance and RAII stay under
two turns and inside the batch noise; RAII's real reading obligation is
not finding the destructor, which sits beside its type, but knowing where
destructors run, an implicit control flow on every exit path, and the
smallness of its measured cost says agents carry that obligation easily.
The reader also holds an oracle the
procedure omits: the compiler. A build confirms a resolution hypothesis
in one turn, the runs used it constantly, and it discounts every
static-resolution step above, though not the derived query, which no
compiler in the repository can check. A derived
query name resolves through a grammar that is in no repository file, so no
number of searches inside the repository ends the read: the reader recalls
the grammar, reads the derivation source priced in category 2, or observes
the runtime, which is the framework row's residence problem. The recorded runs are small enough
to sweep, and their traces show both readings: across all 280 the tool
counts are 911 shell calls, 810 file reads, 412 edits, and 6 writes, the
harness's search tool was never called, and 154 of the shell calls carry a
grep, about two thirds of them searching the source and the rest
filtering file listings, so the agent enumerates and reads at this scale
and still searches when a name is worth finding. File count converts
to reads and reads to turns, which is the mechanical form of the control
result; what stays untested is the price of grep-per-hop reading in a
repository too large to sweep, where the search would be the only
mechanism left. The procedure was written after the first results; recorded in advance
was only the clause that the explicit side's advantage grows with the
number of plausible dispatch targets per call site.

The plain closures fit in a context window beside the work. The framework's
category~3 is text only in reference manuals that are not version-pinned in
any repository, and the training corpus superposes their versions: a coherence
problem, which window growth does not remove, while it does absorb the C++
row's volume (a 1M-token window holds 257k beside the work). That makes the
framework claim the stronger of the two, and the C++ claim a matter of
volume and of counting granularity. The paper's own taxonomy separates them, and the hypothesis in
the abstract rests on the framework row.

\textbf{The web pair, plain JS and React.} The same rule extends to the
browser: the engine (parse, layout, paint) is floor for the reason MySQL
is, and React is counted for the reason Hibernate is, it varies by
version and the application configures it. The subject screen is the
shape System~A uses on every list page: fetch rows, filter, sort,
select, delete, with rendering by a map over the rows joined into
\texttt{innerHTML} and one delegated listener, the idiom of System~A's
own 62,536 lines of JS (387 \texttt{innerHTML} assignments, 118
\texttt{closest()} delegations). Built as behavior-matched twins, equal on the graded projection of
the rendered DOM (count, selection, rows, carets, empty state), the
plain page is 1,092 tokens and the React~19
page is 1,145 in four components, 1,108 in one file: writing concedes
nothing to JSX in this idiom. Behind the React page sit 282,717 tokens
over 29 modules of the runtime's own source, counted on tag v19.1.0
while the runs installed 19.2.8 (the
hooks store, the work loop, reconciliation and key identity, the commit
path, prop application, the delegated synthetic event system, the
scheduler), 2.3 times the Spring closure; the whole runtime is 641,367.
A real third-party subject bounds the twin: the profile screen of a
current RealWorld React implementation (React 19.2, May 2026) closes
over 28 repository files, 16,237 tokens, by mechanical import trace,
and ten of the modules it imports exist in no repository at all: its
API client and schemas are generated at build time from a 1,025-line
OpenAPI specification, category 3 in a new form. The trace is
\texttt{frontend/trace\_subject.py}.

\section{What an agent pays: two bounded experiments}
\label{sec:bounded}

\textbf{Mechanism.} Sixty maintenance runs, in which the agent was told
nothing of the study (Claude Code headless, \texttt{claude-sonnet-5}, four
tasks that name behavior and never a file or a function, five languages,
hidden tests). The source under maintenance was at most 1.1\% of any run's context
(the largest fixture tree, 1,805 o200k tokens, against the smallest run
context), while the median run re-sent 25,603 tokens every
exchange, most of it the fixed preamble the harness sends each time plus the
accumulated conversation, nearly all of it served from prompt cache at a
fraction of full price, so a whole run cost \$0.11. The pilot establishes the
shape of the cost: it is paid per exchange and dominated by what is
carried, and the five representations did not separate beyond
run-to-run spread, though no formal test was run at twelve runs per
language. It cannot establish any claim about scale, because every run sat far
below the window. The cost of a closure beyond the window, retrieval hops and the working
set after retrieval, remains for the experiments named in
Section~\ref{sec:pairs}.

\textbf{Repeated sites.} The don't-repeat-yourself rule, DRY, exists because
a human asked to change $N$ call sites misses some. Twenty-five runs asked for one behavioral change that must
land at every site, at 1, 3, 10 and 30 repetitions plus a factored control,
task text naming no site. Coverage was 30 of 30 in every run that had 30
sites, and 100\% in every condition. This is a ceiling result: at five
runs per condition, a per-run miss probability up to
$\approx$45\% is not excluded (the one-sided 95\% upper bound at 0 of
5). The sites were exact copies one screen apart in one file, and the
task text stated universal scope. The run records sharpen the ceiling further: the 30-site runs
finish in five or six tool calls with zero per-site edits, one
replace-all pass over identical text, so coverage here is a property of
the edit tooling as much as of the model, and site count is irrelevant by
construction. Only one thing stands: the easiest form of the hazard
cannot occur this way; no single pass ends the field form, near-miss
duplicates scattered across files (14 to 53\% normalized against 3.9 to
21.8\% exact in the public corpora), and it is untested.

\section{The shape of the difference, before any agent runs}

Public corpora, 1.74M code lines measured with one script: C carries 28.4
functions per thousand lines, C++ 60.0, Rust 72.5, and the median C++
function has cyclomatic complexity 1: no branch, which is consistent with a
delegation-heavy style though a branchless function can also compute.
Complexity concentrates in C (a maximum of 1,330 in sqlite) and spreads thin
in Rust (25 in tokio). Genre is confounded with language here: the C side is
servers and databases, half the C++ side is header-only template libraries.
One measured contamination is corrected in the shipped data: spdlog vendors
a copy of fmt under \texttt{bundled/}, and vendored trees are excluded from
the corpus counts. The corpus establishes one thing only: the indirection whose cost is in
question is real, large, and ordered from C through C++ to Rust.

System~A shows the low end is an architectural choice, not a property of C:
HTTP to SQL in 3 artifacts, no entity, DTO, mapper, or repository classes,
205 methods returning the response body directly, deployed as 8 jars
totaling 3.3\,MB with a frameworkless 46k-line front end; its deployment
discipline is the build-once-promote model described, with the same system
anonymized the same way, in \cite{promote}. Comparable public
JPA systems declare 73 to 482 Maven artifacts (counted from their POM files
at the recorded commits; one adds 554 npm packages), each a version to track
and a vulnerability feed to watch, none of it visible in a line count.

The author's two Java toolkits from 2009 and 2011, \texttt{aisgedcom} and
\texttt{aisconvert} \cite{aisgedcom,aisconvert} (8,856 code lines, 461
methods between them), carry the same profile in plain Java a decade before
agents existed: median method 7 and 9 lines at complexity 2, complexity
maxima of 41 and 44, no Java file changed since their 2015 import, and both
compile unchanged under javac~21 today (aisgedcom with its legacy source
encoding declared), which is what stable semantics means in practice.

Two guards against over-reading the C rows. Size: the same behavior in this
author's C costs up to 2.7 times his Java in lexical tokens, and the agent
pays that price in full (mincdp: 6,932 o200k against 2,615); compression
shows the excess is mostly repetition (the xz ratio falls to 1.73); and the sites records
(Section~\ref{sec:bounded}) show why exact repetition is safe: one replace-all pass
erases identical copies, and the untested hazard is the near-miss form
no single pass ends. The C advantage claimed here is closure locality, never
size. Runtime: nothing here rests on machine cost.

\section{Why the optimum moved}

Objects are for people: a class is a bundle sized to a human working set,
an interface is a promise a human can hold whole, and an agent reads in
neither unit. Structs are enough for the data and functions over them are
enough for the algorithm, which is how both measured plain systems are
written, down to the persistence layer that returns rows, not objects,
and most visibly in streaming and the hard algorithmics: the ciphers, the
compressors, the codecs are functions over buffers in every codebase,
whatever the language around them. Abstraction layers compensated for specific human limits, and building
them was the correct engineering for the readers of that time; the claim here is that
the optimum moved when the reader changed, not that the layers were a
mistake when they were built. The functions are separable, and they lose their
value separately:

\begin{enumerate}[nosep]
\item \textbf{Reliable mechanical edits} (DRY, central change points). The
sites experiment did not elicit the failure at 30 exact-copy sites; the
near-miss, cross-file form remains open. The indirection cost of every read
remains in either case.
\item \textbf{Compression for small working memory.} Abstraction is a lossy
compressor discovered for the human reader \cite{compressors}; an agent
reads the whole deciding text where it fits beside the work, and
retrieves slices where it does not. The marginal value of hiding detail reaches
zero where the whole truth fits; whether anything changes sharply past
that point is open; what to measure there is the working set after
retrieval, not the window boundary, and the experiments are named in
Section~\ref{sec:pairs}.
\item \textbf{Coordination contracts between teams.} Partially survives:
a module seam can carry the contract a framework boundary carried; how far
team sizes actually fall is not measured here.
\item \textbf{Reuse of hard, stable code.} Survives fully: the platform
floor. The measured plain systems keep abstraction at floors and drop it as
convention.
\end{enumerate}

Corollary, argued rather than measured: languages with stable semantics
(C \cite{knr}, FORTRAN, Ada, plain Java \cite{jls,jdbc}) become viable
again, because everything ever written about them describes the same
language, one coherent body in the training corpus, while a framework's
corpus is a superposition of incompatible versions. The old languages were
expensive for unaided humans, and the expense was exactly the bookkeeping
agents now do. Corpus coherence has no metric in this paper yet; defining
one (how widely the public text about a construct disagrees across versions)
is on the roadmap, and until then this paragraph is argument.
Precedent for the fast death of a compensating abstraction: calculators
ended log tables and slide rules within a decade.

The direction of construction inverts with it: the layers accumulated
because building went bottom-up, each generation standing on the last. With an agent, programming runs
top-down: state the objective, state the guards, and add one bias to the
prompt (``plain C'', ``plain Java with JDBC, streaming'') and the middle
between goal and floor stays thin; the closures of Table~\ref{tab:quad} are
the result of that one-line bias. The bias is necessary: an agent's default is
its training-set average, the industry idiom, and left undirected it
regresses to that mean. Several of the C projects in the author's corpus,
\texttt{mincdp} among them \cite{mincdp}, were rewritten with an AI coder
under the same explicit-style direction: the low-closure style is the
agent's own output once the bias is stated, and the guidance is what
steers it away from the average it learned.

The human work that remains is the objective function (requirements
precise enough to implement) and the guards (what the tests must
assert).
Language choice becomes a measurement: pick what minimizes
tokens-to-trace and corpus incoherence, because nobody is paying the
human learning cost anymore.
Relational access follows the same rule: SQL written against the JDBC
contract \cite{jdbc,date} is one stable text the database runs as written,
which is the plain stacks' zero in categories 2 and 3.

\section{The experiment, run}
\label{sec:pairs}

Six constructs, each a pair: a small C program and the same behavior in
C++ through one construct (a virtual registry, a template, operator
overloads, RAII, inheritance, and a mix of them), plus a control pair
added at review, the same registry classes in a single file. Four tasks
per pair with the same text in both languages, five repetitions, 280
runs, graded by hidden tests the agent never saw; predictions were
written 2026-08-24, before any pair existed. The C++ sides follow the
anchor era's idiom (classes, a registry, \texttt{unique\_ptr}), not
post-C++17 style. The pairs, the grader, and the raw runs of both
executions are beside this paper.

278 of 280 runs passed (the two failures are one C++ cell missing the
same empty-store edge twice), so correctness does not distinguish the sides at
this size. Turn counts do: the C++ side cost more in all seven pairs, six of
six independent constructs, $p = 0.031$, medians 8 against 9, while the
money difference, \$0.105 against \$0.109 a run, does not resolve at
the construct unit.

\begin{table}[H]
\centering
\footnotesize\setlength{\tabcolsep}{3pt}
\begin{tabular}{@{}lrrrr@{}}
\toprule
Construct & files & C & C++ & $\Delta$ \\
\midrule
virtual registry, class per file & 8 & 7.85 & 11.35 & $+$3.50 \\
mixed constructs & 4 & 7.85 & 10.65 & $+$2.80 \\
template in a header & 2 & 6.90 & 8.60 & $+$1.70 \\
RAII & 2 & 8.15 & 9.20 & $+$1.05 \\
inheritance & 2 & 9.15 & 9.90 & $+$0.75 \\
virtual registry, single file (control) & 1 & 6.80 & 7.30 & $+$0.50 \\
operator overloads & 2 & 8.40 & 8.85 & $+$0.45 \\
\bottomrule
\end{tabular}
\caption{Mean turns to termination, 20 runs per construct and language;
all runs but two ended passing; the two failures sit in the single-file
registry's C++ cell and are included in its mean. The C side of both registry rows is the same
program. This execution cost \$29.96 and 117 machine minutes.}
\label{tab:pairs}
\end{table}

Three numbers in the table stand above noise: the registry in eight
files ($+$3.50, $p = 0.003$), the mix ($+$2.80, $p < 0.001$), the
template ($+$1.70, $p = 0.004$). The noise itself is measured: the same
C program, run 20 times twice, differs by 1.05 turns. The control row
answers the section's question: the same registry in one file costs
$+$0.50, indistinguishable from zero, so the agent pays for the walk
across files, not for the virtual dispatch (the cut from 11.35 to 7.30
turns has $p \approx 3\times10^{-4}$). We ran the experiment twice;
three C++ sides were repaired between the executions, and on the
unchanged pairs the small numbers moved (RAII $+$0.50 then $+$1.05), so
we trust the sign, the three large effects, and nothing finer. One bias
ran toward the null and survived: the scattered idiom's file names hint
where to look.

The pairs also give the behavior-matched closure comparison the four
systems cannot: over the seven pairs the C sides total 3,898 o200k tokens
of deciding text against 4,547 for the C++ (source files only, build
files excluded), a factor of 1.17 with the single-file control included, 1.13 over the
six idiomatic pairs, and in two pairs the C++ side is the smaller
(template 0.82, inheritance 0.96). Held to identical behavior, the idiom
costs 13 to 17\% more text; the 2.8 times
between the whole tools is therefore mostly feature scope and what a
style accretes at system size, not the constructs.

Of the recorded predictions, only the cost direction held: the prediction said C++ would be cheaper on
local tasks, and it was not; the predicted failures never came. A caution on scale: every run sat far below the context window,
and an agent retrieves slices rather than loading closures, so what a
large closure costs is untested here. The volume experiment that would
test it, cross-module tasks on taskwarrior itself against its plain
twin, remains open.

\subsection{The endpoint, two stacks}
\label{sec:endpoint}

The second named experiment ran: one endpoint of the System~A kind, an
item list and an item create behind bearer-token authentication over a
local SQL store, built twice with the same routes, schema, seed rows and
observable behavior. The plain side is one file, JDK
\texttt{HttpServer} and JDBC, 108 lines. The framework side is the
stack of Section~\ref{sec:closures} (Spring Boot~3.3 with Spring Data JPA; the closure
there was measured on Boot~4), seven
classes in the current idiom (a \texttt{OncePerRequestFilter} for the
token, a repository interface with derived queries, a
\texttt{@RestController}). Equivalence at the HTTP boundary, status,
parsed body and the pinned header, was verified over eight probe cases
of the contract before any run (\texttt{equiv.py}); outside the
contract the sides diverge, a POST body missing a field getting 400
from plain's visible validation and 500 from Spring's database
constraint, a case no task touches. Four maintenance tasks,
the same text on both sides: add a field to the response, add a
combinable query-string filter, change the authentication failure
contract, and find an injected prefix-match defect. Five repetitions,
40 runs, one model (claude-sonnet-5, headless), grading by hidden HTTP
checks against a live server the agent never saw; the whole experiment
then ran a second time under the same harness five days later, another
40 runs. The criterion was recorded before the build: the hypothesis
fails if the framework tasks do not cost more than the plain ones.

\begin{table}[H]
\centering
\footnotesize\setlength{\tabcolsep}{3.5pt}
\begin{tabular}{lccccc}
\toprule
 & \multicolumn{2}{c}{batch 1} & \multicolumn{2}{c}{batch 2} & \\
Task & plain & Spring & plain & Spring & pass \\
\midrule
add a response field & 6.6 & 8.4 & 10.0 & 7.4 & 10/10\;9/10 \\
add a filter & 11.0 & 12.0 & 9.6 & 15.2 & 10/10\;6/10 \\
change the 401 contract & 6.0 & 4.2 & 9.2 & 4.2 & 10/10\;10/10 \\
find the injected defect & 5.0 & 9.0 & 7.6 & 7.8 & 10/10\;10/10 \\
\bottomrule
\end{tabular}
\caption{Mean turns per task, five repetitions per batch. Pass columns:
plain then Spring, both batches pooled.}
\label{tab:endpoint}
\end{table}

The turn count did not survive the second batch. In the first the
framework side cost 1.25 turns more per run (7.15 against 8.40) and
5.3 seconds more wall time, both above the within-batch noise with
task as the designed blocking factor ($p = 0.023$ turns, $p = 0.022$
wall, one-sided permutation, 100{,}000 resamples). The second batch
reversed the sign, $-0.45$ turns ($p = 0.62$), and the 80 runs
combined give $+0.40$ ($p = 0.31$). The reversal is what the pairs
experiment's own calibration predicts: its control twin, the same C
program in two cells of 20, differs from itself by 1.05 turns, and the
same machinery declares that null difference significant, so a
within-batch $p$ measures exchangeability inside one batch, not
stability across batches. The recorded criterion is met in the first
batch and not in the second: the cost claim is void. Cumulative
context ran 21 thousand tokens higher on the framework side in the
first batch and does not resolve ($p = 0.13$).

Correctness separates the sides in count as well as in kind: 35 of 40
against 40 of 40, all five failures on one side (one-sided
hypergeometric $p = 0.027$), every one silent, every one a wire name
derived from a Java identifier by a convention no application file
states. Four are the same miss on the filter task, two per batch: the
agent added \texttt{@RequestParam(required =
false) Integer minLen} and correct JPQL behind it, but Spring binds a
query parameter by the Java parameter name (this build, like every
Boot build, compiles with \texttt{-parameters}; without that flag the
same code fails loudly at the request instead), the URL says
\texttt{min\_len}, so the
parameter stayed null and the endpoint returned the unfiltered list
with status 200. The passing repetitions wrote the wire name
explicitly (\texttt{name =} or \texttt{value = "min\_len"}). The fifth
is the same convention on the response side: asked for a
\texttt{name\_length} field, the agent added
\texttt{getNameLength()} and Jackson's bean naming emitted
\texttt{nameLength}, field present, name wrong, status 200, while its
four passing siblings wrote \texttt{@JsonProperty("name\_length")}
explicitly. Nothing in the
compiler, the runtime or the response flagged the five that did not:
the deciding rules are category-3 conventions of Section~\ref{sec:measure},
and this is what their failure looks like, silent wrong rows. The plain
side parses its query string and writes its JSON keys in visible code,
so the same mistakes had no place to hide, ten of ten.

The convention residence is itself counted for this stack: the fifteen
conventions on the endpoint's path, both naming rules among them, are
stated in \textbf{31,535 o200k tokens} of documentation, measured from
the documentation sources at the exact BOM-resolved tags (Framework
6.1.13, Data JPA and Commons 3.3.4, Boot 3.3.4, Hibernate 6.5.2,
Jackson 2.17.2; whole pages, or the section stating the rule, so a
lower bound; H2 is floor as MySQL is elsewhere; \texttt{cat3.py} /
\texttt{cat3.json} beside this paper). That is four times the plain
side's whole deciding text, and the page stating the
\texttt{-parameters} rule is among the counted rows.

The direction is not uniform and the exception is instructive: on the
401-contract task Spring was cheaper in both batches (4.2 against 6.0
turns, then 4.2 against 9.2), because
the framework's filter puts the whole authentication contract in one
small class while the plain side's is inline in a longer file. That
localization win is the only per-task turn effect that repeated; the
defect hunt's $+4.0$ of the first batch fell to $+0.2$ in the second.
Where the framework localizes a change it can win the turn count;
where the behavior crosses its naming conventions it fails silently.
Two scope cautions: both projects are small, so
the 51 to 124 thousand tokens of category-2 text from Section~\ref{sec:closures} never
had to be read, and Spring is in every training corpus, so these runs
price recalling the framework's conventions, not reading them; the
wall gap also folds in the slower build the framework brings
(Gradle against \texttt{javac}) and the harness ran three runs in
parallel, which adds contention to wall time on both sides. The build
scripts, tasks, grader, equivalence probes and
raw runs are in \texttt{endpoint/} beside this paper.

\subsection{The web pair, run}
\label{sec:web}

The third pair ran twice; the predictions were written before the
twins, the same day, and entered the repository with them. First the
twins: four maintenance tasks mirroring the endpoint four, five
repetitions, 40 runs, graded by hidden scenarios reading the rendered
DOM in headless Chrome. All 40 passed. Turns did not separate, 6.95
against 6.85 ($p = 0.91$), and React won the add-a-field task, 5.6
against 7.6 turns: the component boundary puts that change in one
small file.

\begin{figure}[!t]
\footnotesize
\begin{verbatim}
plain JS             React
1.1k, all in page    1.1k + 283k runtime

event listener       {synthetic event}
  |                    |
handler              {priority lane}
  |                    |
filter + sort        state setter
  |                    |
template string      {hook state store}
  |                    |
innerHTML write      component function
  |                    |
browser paints       filter + sort
                       |
                     {reconciler diff}
                       |
                     {key identity}
                       |
                     {commit phase}
                       |
                     {prop application}
                       |
                     browser paints

  |   a step to a name at the call site,
      resolvable by one search
 {}  a step resolvable only in the
     framework's source or conventions
\end{verbatim}
\caption{The third pair, the same interaction on the web: a keystroke in
the filter until the browser paints the rows. The browser engine (parse,
layout, paint) is floor on both sides by the rule that makes MySQL floor
in the Java pair. Every plain step is in the page's own 1,092 tokens; the
braced steps resolve in the 282,717 tokens of the runtime's source
(react, react-dom, the reconciler, the scheduler, at the measured
version), or, for the delegated event's target and the controlled
input's value tracking, in version convention. The two sides' own code
is the same size (1,092 against 1,145 tokens, matched behavior,
identical rendered DOM); the difference is the machinery behind the
braces.}
\label{fig:webpaths}
\end{figure}

Then the version conventions directly: two tasks whose correct idiom
changed between React 18 and 19 (\texttt{defaultProps} is silently
ignored on 19; a ref passed as a plain prop is silently stripped on
18), each against both pinned versions (18.3.1, 19.2.8) plus the plain
control, 30 runs. All passed, zero version-wrong choices: the model's
habitual idioms are the version-portable ones, the inline default and
\texttt{forwardRef}, and only two of the ten version-critical runs
looked at \texttt{package.json}. What separated is the boundary:
focusing a child component's input from the parent cost 12 and 13 mean
turns against plain's uniform 6, 32 and 36 seconds against 18, 323 and
353 thousand cumulative context tokens against 181, and \$0.14 a run
against \$0.09, every run passing. The web pair repeats the pairs
result, not the endpoint one: nothing broke, and the agent pays for
the walk across the boundary. The twins, the grader, the equivalence
probe and the subject trace are in \texttt{frontend/}; the run
records of all experiments ship with this paper's Zenodo deposit
(\texttt{run\_records.zip}).

\section{Related measurements}

The nearest prior work measures token consumption during runs. A study of
agent trajectories across Java, OCaml, Python and Rust on LiveCodeBench
problems \cite{tokenmax} measures consumed tokens per language, difficulty
controlled, and finds that the language with the longest programs is not
the one that costs the most, which agrees with our pairs result. An
analysis of token consumption on SWE-bench Verified \cite{spend} finds JVM
projects dearer than Python largely through verbose build and dependency
logs, an ecosystem cost adjacent to our category 3. Cross-language
benchmarks such as MultiPL-E \cite{multiple} compare pass rates, not cost;
SWE-bench \cite{swebench} prices agents on real repositories with the
language mostly fixed; and the token-cost question was asked for natural
languages through tokenizers \cite{petrov}. None of these measures the
deciding text of a real endpoint before any run, or classifies it by
residence, which is what tokens-to-trace adds. The two kinds of count are
complementary, and the static one exists for a reason the spend cannot
serve: consumed tokens are a property of the model, the harness and the
date, and in our own pilot, whose fixtures are small, they were mostly
harness, the source never above 1.1\% of a context, though at
repository scale the reading itself dominates the spend; the deciding text is a property of the code
alone, computable before anything is built. The difference is the
context a path demands: for the plain C path, six files the
agent can read whole; for the same logic on Spring, the project's part,
the framework's interpreters, and conventions no context can hold. A
spend count cannot see that difference, because an agent answering from
its weights makes the missing context look free.

\section{Threats to validity}

The C++ we measured is the idiom as it is practiced \cite{stroustrup},
not the language itself; C-style code compiled as C++ would get the C
row's closure. The public subjects show the idiom's average, while both
plain systems were written to a stated style that deliberately leaves
that average, so the comparison is between what teams write by default
and what one line of direction produces. Anyone can state the same line.

One rater derived the closures, and he is the author of ais and one of
System~A's developers. Each system contributes one endpoint, and he chose
which. To check him, a second rater re-derived two closures from the
definition alone, without seeing the manifests. On the C++ file list the
two agree on 97 of 103 files: the second rater added the status column
and the JSON module, and dropped the dependency graph and one
implementation file. On the framework path the second rater found 50
classes where the paper counts 27, with 25 of the 27 among them, and his
50 measure 373k tokens against the printed 124k. The printed category 2
is therefore the conservative trace, and its being an undercount is
confirmed. taskwarrior is a single registry-heavy codebase, pinned at
2.6.2 because that is the last all-C++ release; leveldb would have
measured smaller. On the plain side the own-code objection is bounded by
the third-party row: Monocypher's encrypt path, written by another
author, closes at the same few thousand tokens as the author's two
systems.

Where the platform floor sits is a judgment, so both ends are printed:
moving it changes ais by 2.5 times (5.8k to 14.4k) and Spring's category
2 by 2.4 times (124k to 51k). A stricter floor would also turn the plain
zeros into small numbers, because \texttt{sql\_mode}, \texttt{getopt}
and the build flags were excluded everywhere, and that omission is
largest for C. SQL itself is an interpreted notation like the
derived-query grammar; it is excused here only because its meaning has
been stable for decades, and stability is the property the taxonomy
wants; residence is its measurable proxy. Category 2 counts
whole classes from one trace that no coverage run has verified, so it
bounds in neither direction, and the C++ row is whole files, with the
same-granularity comparator printed beside it.

The measure prices reading and nothing else. It does not price the
knowledge a reader brings from outside: the years a person spends
becoming proficient in C++ or in a framework, and the training corpus an
agent holds for the same purpose, sit on top of every number here. The
size of that implicit context can itself be compared by the defining
books, K\&R's 272 pages against Stroustrup's 1,368 \cite{knr,stroustrup},
and a framework's reference has no fixed size at all: it moves with the
version. It
does not price security, which is its own dimension: framework defaults
replace hand-rolled code whose defect classes are real, and System~A
answers only part of that by keeping its secrets in one place, outside
the repository and the artifact. It cannot credit what a type system
removes: \texttt{const} and RAII settle questions the reader of C must
verify by hand, so the C++ closures are overcounted and the manual
bookkeeping in the ais engine is undercounted. And in certification work
under DO-178C or ISO~26262, every duplicated site still has to be
verified separately, whatever the agent's edit coverage.

Both agent experiments of Section~\ref{sec:bounded} sat far below the context window and
passed everything, so they bound mechanisms and cannot separate
representations. In the pairs the C++ file count changes together with
the construct, one class per file being the registry idiom, so the top
of Table~\ref{tab:pairs} prices construct and layout together. Spring
and its idiom are in every training corpus, so the endpoint runs price
recalling the framework's conventions rather than reading them, and
recall is where its five failures happened. The web twins are
author-written and small, the real subject bounds only the closure, and
the stale-recall condition tested two version boundaries at five
repetitions per cell, so its zero is a bound, not a law; the DOM
projection the web grader compares is a choice, and a different
projection could grade differently; and no plain cell with the same
indirection, the input extracted into its own module, was run, so the
boundary's price is measured against the unfactored plain page only. One model ran, on one date;
agent cost is a property of the model and the representation together,
and the web result says the sharpest version of that: the version traps
missed because this model's habits are version-portable, and another
model's habits need not be.

\section{Conclusion and recommendation}

It is shown, in counts anyone can recompute, that writing in structs and
functions, in an old language with stable meaning, is far more compact
in the agent's memory: the whole deciding text of a path costs a few
thousand tokens and leaves the window to the work, while the layered styles spend
the window itself (257k) or keep the deciding text where no repository
holds it. The experiments put small numbers on the syntax: held to identical
behavior, the C++ sides carry only 1.17 times the C sides' text; below
the window a construct of indirection kept in one or two files costs the
agent under two turns, differences under a turn sit inside run-to-run
noise, almost nothing breaks, and the one large cost,
measured and controlled, is the file scatter the idiom brings. On the
Java stacks the endpoint ran in two batches of 40; the turn cost
reversed its sign between them and is not a finding, while the
failures replicated: five silent wrong answers over 80 runs, all the
framework side's, each a wire name derived from a Java identifier by a
convention no source file states, against zero on the plain side
($p = 0.027$); the one task its filter class localizes it won in both
batches, and the asymmetry is in Section~\ref{sec:endpoint}. On the web
pair nothing broke and nothing separated but the boundary, twice the
turns to cross a component (Section~\ref{sec:web}), so the third pair
joins the pairs result, not the endpoint one. What a closure that outgrows the window costs is still open, and its scope is
stated: an agent retrieves slices rather than loading wholes,
so the question at scale is the working set after retrieval and the text
no retrieval reaches, and the volume experiment that decides it is named
above.

The recommendation, scoped to read cost and to teams working with agents:
state the bias in the prompt (``plain C'', ``plain Java with JDBC,
streaming'', ``plain vanilla JavaScript''), keep abstraction only where a platform floor earns it, spend the human
effort on the objective (the requirements) and the guards (what the tests
must assert; the working practices this compresses are published as
recipes \cite{recipes}), and when choosing a stack, measure the closure of one
representative path before arguing about taste.

\section*{Method note}

Closures measured 2026-08-25: System~A private tree; \texttt{ais} at head
\cite{ais}; taskwarrior 2.6.2; the Spring application at v2.1.1 with
framework source from the published sources jars of the BOM-resolved
versions; every manifest records its version or commit, the public subjects
are re-fetchable from the cited repositories, and the exact checkouts are
available from the author on request. Tokens are o200k via
tiktoken; manifests with per-row lines, characters, and tokens are
\texttt{closures.py} / \texttt{closures.json}. Corpus, duplication, JPA,
pilot and site-coverage numbers come from the scripts and JSON beside this
paper (\texttt{measure.py}, \texttt{dup.py}, \texttt{external.py},
\texttt{jpa.py}, \texttt{pilot/}, \texttt{sites/}), each recording its
commits and raw runs; the endpoint runs (claude-sonnet-5, two batches
of 40: 2026-08-25 and 2026-08-30) and both stacks' sources are
\texttt{endpoint/}; the category-3 documentation count is
\texttt{cat3.py} / \texttt{cat3.json}, fetching the documentation
sources at the BOM-resolved tags; the web pair
(claude-sonnet-5, 2026-08-26; React 19.2.8 and 18.3.1 via vite, chrome
headless grading; closure counted on the react source tag v19.1.0 and
on yurisldk/realworld-react-fsd at 969709a) is \texttt{frontend/}; the blind second rater's derivation is
\texttt{second\_rater.txt}; the Maven-artifact counts are from the projects' POMs
and not yet scripted. Region boundaries are function, block, or
DDL-statement granularity; whole files where stated.

\begin{thebibliography}{99}
\bibitem{knr} B.~Kernighan, D.~Ritchie. \emph{The C Programming Language},
2nd ed. Prentice Hall, 1988.
\bibitem{robbins} A.~Robbins. \emph{Linux Programming by Example: The
Fundamentals}. Prentice Hall, 2004.
\bibitem{stroustrup} B.~Stroustrup. \emph{The C++ Programming Language},
4th ed. Addison-Wesley, 2013.
\bibitem{misra} MISRA. \emph{MISRA C:2012 Guidelines for the use of the C
language in critical systems}. HORIBA MIRA, 2013.
\bibitem{jls} J.~Gosling, B.~Joy, G.~Steele, G.~Bracha. \emph{The Java
Language Specification}, 2nd ed. Addison-Wesley, 2000.
\bibitem{jdbc} S.~White, M.~Hapner. \emph{JDBC 2.1 API Specification}. Sun
Microsystems, 1999.
\bibitem{date} H.~F.~Korth, A.~Silberschatz. \emph{Database System
Concepts}. McGraw-Hill, 1986.
\bibitem{monocypher} L.~Vaillant. \emph{Monocypher} 4.0.3, a
cryptographic library in C99. \url{https://monocypher.org}.
\bibitem{taskwarrior} \texttt{taskwarrior}, an open-source command-line
task manager, in development since 2006; 2.6.2 is the last all-C++ release.
\url{https://github.com/GothenburgBitFactory/taskwarrior}.
\bibitem{realworld} \texttt{realworld-springboot-java} v2.1.1, a Spring
Boot implementation of the RealWorld specification (one Medium-style
application, implemented identically across many stacks).
\url{https://github.com/raeperd/realworld-springboot-java}; the
specification: \url{https://github.com/gothinkster/realworld}.
\bibitem{tokenmax} \emph{The Best Programming Language for Tokenmaxxing:
An Investigation of Coding Agent Behavior Across Programming Languages}.
arXiv:2607.22807, 2026.
\bibitem{spend} \emph{How Do AI Agents Spend Your Money? Analyzing and
Predicting Token Consumption in Agentic Coding Tasks}. arXiv:2604.22750,
2026.
\bibitem{multiple} F.~Cassano et al. \emph{MultiPL-E: A Scalable and
Polyglot Approach to Benchmarking Neural Code Generation}. IEEE TSE, 2023.
\bibitem{swebench} C.~Jimenez et al. \emph{SWE-bench: Can Language Models
Resolve Real-World GitHub Issues?} ICLR, 2024.
\bibitem{petrov} A.~Petrov et al. \emph{Language Model Tokenizers
Introduce Unfairness Between Languages}. NeurIPS, 2023.
\bibitem{tiktoken} tiktoken, OpenAI's open-source tokenizer library;
\texttt{o200k\_base} is the encoding of OpenAI's current models.
\url{https://github.com/openai/tiktoken}.
\bibitem{aisstyle} \emph{AIS: Coding Style and Ideology}. In \cite{ais},
\texttt{doc/dev/STYLE.md}.
\bibitem{ais} \texttt{ais}, an associative index in C99.
\url{https://github.com/Anode1/ais}. The five-language bench and its
\texttt{LANG\_COMPARISON} note are under \texttt{tests/perf}.
\bibitem{mincdp} \texttt{mincdp}, a Chrome DevTools Protocol client, one
file per language. \url{https://github.com/Anode1/mincdp}.
\bibitem{aisgedcom} \texttt{aisgedcom}, a GEDCOM genealogy toolkit in Java,
2009--2011. \url{https://github.com/Anode1/aisgedcom}.
\bibitem{aisconvert} \texttt{aisconvert}, a personal genomics toolkit in
Java, 2009. \url{https://github.com/Anode1/aisconvert}.
\bibitem{recipes} \texttt{agent-recipes}: working practices for coding
with AI agents, generalized from the systems measured here.
\url{https://github.com/Anode1/agent-recipes}.
\bibitem{compressors} V.~Gavrilov. \emph{Intelligence Is the Discovery of
Compressors}. Zenodo, 2026. \url{https://doi.org/10.5281/zenodo.20440110}.
\bibitem{promote} V.~Gavrilov. \emph{The Artifact Promotion Control Model:
An Implementation Case Study}. Zenodo, 2026.
\url{https://doi.org/10.5281/zenodo.20528903}.
\end{thebibliography}
\end{document}
